%! TeX program = lualatex
\documentclass{tesis}
\usepackage{fontspec}
\setmainfont{Times New Roman}

\newcommand{\autor}{Jeniffer Paulina Yaguana Caraguay}
\newcommand{\cedulaAutor}{1106113697}
\newcommand{\tutor}{Leonardo David Torres Valverde}
\newcommand{\lugarFechaPrelims}{Ambato, julio 2024}
\newcommand{\tema}{Prototipo de un Sistema de Monitoreo y Control para la gestión de consumo de agua en los cultivos de plantas del Vivero Michita mediante Lógica Difusa}
\newcommand{\elaboradopor}{Elaborado por: el investigador}


\begin{document}

%Start counting pages from 2 in roman numbers
\pagenumbering{roman}

%A. Paginas preliminares
\addcontentsline{toc}{chapter}{\bfseries \uppercase{PORTADA}}
\begin{titlepage}
	\begin{center}

		\includegraphics[width=50mm]{resources/images/logo-uta.png}

		\textbf{\uppercase{
				Universidad Técnica de Ambato \\
				\vspace{5mm}
				Facultad de Ingeniería en sistemas, electrónica e industrial \\
				\vspace{5mm}
				carrera de software \\
			}}

		\vspace{5mm}
		\textbf{Tema:}
		\vspace{5mm}

		\begin{tabular}[c]{>{\centering\arraybackslash}p{0.97\textwidth}}
			\toprule
			{\bfseries\MakeUppercase\tema} \\
			\bottomrule
		\end{tabular}

    \vspace{24pt}
		Trabajo de titulación modalidad Proyecto de investigación, presentado
		previo a la obtención del título de Ingeniero de Software

		\vspace{5mm}
		\textbf{\uppercase{área:}} Software\\
		\vspace{5mm}
		\textbf{\uppercase{línea de investigacion:}} Desarrollo de Software\\
		\vspace{5mm}
		\textbf{\uppercase{autor:}} \autor \\
		\vspace{5mm}
		\textbf{\uppercase{tutor:}} Ing. \tutor, Mg.\\

		\vfill
		Ambato - Ecuador\\
		febrero - 2024

	\end{center}
\end{titlepage}


\setcounter{page}{2}
% %TODO
\chapter*{aprobación del tutor}
\addcontentsline{toc}{chapter}{\bfseries \uppercase{aprobación del tutor}}

En calidad de tutor del trabajo de titulación con el tema:
{\MakeUppercase\tema},
desarrollado bajo
la modalidad Proyecto de Investigación, por el señor \autor,
estudiante de la Carrera de Software, de la Facultad de Ingeniería en Sistemas, Electrónica e Industrial, de la
Universidad Técnica de Ambato, me permito indicar que el estudiante ha sido tutorado
durante todo el desarrollo del trabajo hasta su conclusión, de acuerdo a lo dispuesto en
el Artículo 17 del Reglamento para la Titulación de Grado en la Universidad Técnica
de Ambato y el numeral 6.3 del instructivo del reglamento referido.

\begin{flushright}
	\lugarFechaPrelims.
\end{flushright}

\vspace*{5cm}
\begin{center}
	\begin{tabular}{c}
		\hbox to 7cm{\leaders\hbox to 5pt{\hss - \hss}\hfil} \\
		Ing. \tutor, Mg.\\
		TUTOR
	\end{tabular}
\end{center}

% %TODO
\chapter*{autoría}
\addcontentsline{toc}{chapter}{\bfseries \uppercase{autoría}}
El presente trabajo de titulaci\'on con el tema: {\MakeUppercase\tema} es absolutamente original, auténtico y personal y ha observado
los preceptos establecidos en la Disposición General Quinta del Reglamento para la
Titulación de Grado en la Universidad Técnica de Ambato. En tal virtud, el contenido,
efectos legales y académicos que se desprenden del mismo son de exclusiva
responsabilidad del autor.\begin{flushright}
	\lugarFechaPrelims.
\end{flushright}

\vspace*{5cm}
\begin{center}
	\begin{tabular}{c}
		\hbox to 7cm{\leaders\hbox to 5pt{\hss - \hss}\hfil} \\
		Richard Manuel Carrión Valarezo                      \\
		C.C. 0705023406                                      \\
		AUTOR
	\end{tabular}
\end{center}

% %TODO
\chapter*{derechos de autor}
\addcontentsline{toc}{chapter}{\bfseries \uppercase{derechos de autor}}
Autorizo a la Universidad Técnica de Ambato para que reproduzca total o parcialmente
este trabajo de titulación dentro de las regulaciones legales e institucionales
correspondientes. Además, cedo todos mis derechos de autor a favor de la institución
con el propósito de su difusión pública, por lo tanto, autorizo su publicación en el
repositorio virtual institucional como un documento disponible para la lectura y uso
con fines académicos e investigativos de acuerdo con la Disposición General Cuarta
del Reglamento para la Titulación de Grado en la Universidad Técnica de Ambato.
\begin{flushright}
	\lugarFechaPrelims.
\end{flushright}

\vspace*{5cm}
\begin{center}
	\begin{tabular}{c}
		\hbox to 7cm{\leaders\hbox to 5pt{\hss - \hss}\hfil} \\
		Richard Manuel Carrión Valarezo                      \\
		C.C. 0705023406                                      \\
		AUTOR
	\end{tabular}
\end{center}

% %TODO
\chapter*{aprobación del tribunal de grado}
\addcontentsline{toc}{chapter}{\bfseries \uppercase{aprobación del tribunal de grado}}
En calidad de par calificador del informe final del trabajo de titulación presentado por el señor \autor, estudiante de la Carrera de Software,
de la Facultad de Ingeniería en Sistemas,
Electrónica e Industrial, bajo la Modalidad Proyecto de Investigación, titulado
	{\MakeUppercase\tema}
nos permitimos informar que el trabajo ha sido revisado
y calificado de acuerdo al Artículo 19 del Reglamento para la Titulación de Grado en
la Universidad Técnica de Ambato y el numeral 6.4 del instructivo del reglamento
referido. Para cuya constancia suscribimos, conjuntamente con la señora Presidente
del Tribunal.
\begin{flushright}
	\lugarFechaPrelims.
\end{flushright}


\vspace*{4cm}
\begin{center}
	\begin{tabular}{c}
		\hbox to 7cm{\leaders\hbox to 5pt{\hss - \hss}\hfil} \\
		Ing. Elsa Pilar Urrutia Urrutia, Mg.                 \\
		PRESIDENTE DEL TRIBUNAL                              \\
	\end{tabular}
\end{center}
\vspace*{20mm}
\begin{center}
	\begin{tabular}{c c}
		\hbox to 7cm{\leaders\hbox to 5pt{\hss - \hss}\hfil} & \hbox to 7cm{\leaders\hbox to 5pt{\hss - \hss}\hfil} \\
		Ing. Carlos Israel Nuñez Miranda, Mg.                & Ing. Leonardo David Torres Valverde, Mg.             \\
		PROFESOR CALIFICADOR                                 & PROFESOR CALIFICADOR
	\end{tabular}
\end{center}

% \newpage
\vspace*{\fill}
\begin{flushright}
  \footnotesize
	\begin{minipage}{0.5\textwidth}
		\begin{flushright}
      \phantomsection
			\uppercase{\textbf{Dedicatoria}}
		\end{flushright}
		\itshape

		\bigbreak
		\begin{flushright}
			\textbf{Richard Manuel Carrión Valarezo}
		\end{flushright}
	\end{minipage}
\end{flushright}
\vspace*{\fill}
\addcontentsline{toc}{chapter}{\bfseries DEDICATORIA}

% \newpage
\vspace*{\fill}
\begin{flushright}
	\footnotesize
	\begin{minipage}{0.5\textwidth}
		\begin{flushright}
			\phantomsection
			\uppercase{\textbf{Agradecimientos}}
		\end{flushright}
		\itshape


		\bigbreak
		\begin{flushright}
			\textbf{Richard Manuel Carrión Valarezo}
		\end{flushright}
	\end{minipage}
\end{flushright}
\vspace*{\fill}
\addcontentsline{toc}{chapter}{\bfseries AGRADECIMIENTO}


%Titulo de la tabla de contendios centrado, bold y mayusculas
\newpage
\tableofcontents
\addcontentsline{toc}{chapter}{\bfseries\uppercase{índice general de contenidos}}

\let\origaddvspace\addvspace
\newpage
\renewcommand{\addvspace}[1]{}
\listoftables
\renewcommand{\addvspace}[1]{\origaddvspace{#1}}
\addcontentsline{toc}{chapter}{\bfseries\uppercase{índice de tablas}}
\newpage
\renewcommand{\addvspace}[1]{}
\listoffigures
\renewcommand{\addvspace}[1]{\origaddvspace{#1}}
\addcontentsline{toc}{chapter}{\bfseries\uppercase{índice de figuras}}
% \newpage
% \listofappendices
% \addcontentsline{toc}{chapter}{\bfseries\uppercase{índice de anexos}}

\newpage
\listofappendixs
\addcontentsline{toc}{chapter}{\bfseries\uppercase{índice de anexos}}

\newpage
\renewcommand{\addvspace}[1]{}
\lstlistoflistings
\renewcommand{\addvspace}[1]{\origaddvspace{#1}}
\addcontentsline{toc}{chapter}{\bfseries\uppercase{índice de códigos}}

%Numeros arabicos despues de los prelims
% \chapter*{resumen ejecutivo}
\addcontentsline{toc}{chapter}{\bfseries\uppercase{resumen \uppercase{ejecutivo}}}
Las Unidades de Producción de la Universidad Técnica de Ambato son entidades
autogestionadas que generan ingresos a partir de la oferta de productos y servicios, 
representando un beneficio significativo para la universidad. Sin embargo, enfrentan desafíos en la
gestión de la información y carecen de
un medio que les permita compartir información con la dirección financiera, encargada de
recibir los pagos, para su visualización y verificación, lo cual supone un problema a la
hora de completar el proceso de entrega de productos y servicios.
\bigbreak
El presente proyecto propone la implementación de una aplicación web que establezca un 
proceso en línea para adquirir y procesar los pagos de servicios y productos, 
permitiendo a los usuarios adjuntar y verificar información relativa a los 
procesos que les corresponden. Además, se busca agilizar el registro y 
validación de datos mediante el reconocimiento óptico de caracteres (OCR) y 
procesamiento de imágenes.
\bigbreak
La aplicación web se divide en dos componentes principales: backend y frontend. En el 
backend, se utiliza Python con FastAPI para el desarrollo de la interfaz de programación de aplicaciones 
(API), integrando OCR con Pytesseract y OpenCV. 
Adicionalmente se empleó PostgreSQL como base de datos y MinIO para gestionar imágenes y archivos. 
En el frontend, se emplea React y bibliotecas adicionales.
\bigbreak
Finalmente, como resulatado se obtuvo una aplicación web que minimiza los incovenientes del proceso
actual y que, adicionalmente, descarta el uso de documentación física evitando así posibles pérdidas de
información que demoren la culminación del proceso.
\vfill
\textbf{Palabras clave:} Rappid Application Development, OCR, Tesseract, procesamiento de imagenes, aplicación web

% \newpage
\chapter*{abstract}
\addcontentsline{toc}{chapter}{\bfseries \uppercase{ABSTRACT}}
The Unidades de Producción of the Universidad Técnica de Ambato  are self-managed 
entities that generate income through the offering of products and services, 
representing a significant benefit for the university. However, they face challenges in 
information management and lack a means to share information with the finance 
department, responsible for receiving payments, for visualization and verification. This 
poses a problem when completing the process of delivering products and services.
\bigbreak
The current project proposes the implementation of a web application to establish an 
online process for acquiring and processing payments for services and products. This 
allows users to attach and verify information related to their respective processes. 
Additionally, the project aims to expedite the registration and validation of data 
through optical character recognition (OCR) and image processing.
\bigbreak
The current project proposes the implementation of a web application to establish an 
The web application is divided into two main components: backend and frontend. In the 
backend, Python with FastAPI is used to develop the Application Programming Interface 
(API), integrating OCR with Pytesseract and OpenCV. Additionally, PostgreSQL is employed 
as the database, and MinIO is used for managing images and files. In the frontend, React 
and additional libraries are utilized.
\bigbreak
In conclusion, the result is a web application that minimizes the drawbacks of the 
current process and, additionally, eliminates the need for physical documentation, 
thereby avoiding potential information losses that could delay the completion of the 
process.
\vfill
\textbf{Keywords:} Rappid Application Development, OCR, Tesseract, image processing, web application



%B. Contenidos
%Change counting to arabic numbers
\newpage
\renewcommand{\thepage}{\arabic{page}}% Roman page numbers
\setcounter{page}{1}

\chapter{MARCO TEÓRICO}

\section{Tema de investigación}
 {\MakeUppercase\tema}

\subsection{Planteamiento del problema}
El agua desempeña un papel fundamental en la agricultura. Dicho esto, la agricultura es uno de los principales consumidores de este recurso hídrico. En \cite{noauthor_gestion_2018} menciona la Organización de las Naciones Unidas para la Alimentación y la Agricultura (FAO) que la actividad agrícola emplea aproximadamente el 70\% del suministro mundial de agua dulce. Es por ello, que la cantidad de agua utilizada en los campos de cultivo varía según diversos factores, como el tipo de cultivo, las condiciones climáticas,
la calidad del suelo y el método de riego empleado. Por lo tanto, la FAO \cite{noauthor_gestion_2018} indica que para aumentar la eficiencia en el uso del agua, es necesario reducir su uso entre el 25\% y el 40\% de este recurso.
\bigbreak
En \cite{temperatura_nodate} señala que las condiciones climáticas del Ecuador varían significativamente entre regiones y estaciones. En este sentido, la diversidad climática influye directamente en la disponibilidad de agua, afectando así el crecimiento adecuado de los cultivos vegetales en viveros. Es por eso que en el artículo \cite{c_estudio_2018}, se menciona que en Mulaló el 83\% de los agricultores están haciendo un uso excesivo del agua de riego, debido a que disponen de cantidades superiores a las necesarias tanto para los cultivos como para la extensión de tierra irrigada.

\bigbreak
La gestión ineficiente del agua puede disminuir la productividad agrícola en viveros. De este modo, cuando los cultivos no reciben la cantidad necesaria de agua, se reduce la eficiencia en la producción y se compromete el desarrollo deficiente de las plantas ornamentales, forestales y frutales. Por otro lado, en \cite{luz_cultivo_2014} resalta que el exceso de agua puede provocar asfixiar a las plantas y generar pudriciones en las raíces.

\bigbreak
La gestión del riego de agua en el vivero Michita se realiza de forma manual, lo que impide un control del consumo de agua por cada tipo de planta. Además, la falta de información  sobre la cantidad de agua requerida para los diferentes tipos de plantas dificulta la gestión eficiente de este recurso hídrico. 

\section{Antecedentes investigativos}
% Luego de un análisis en repositorios digitales, plataformas académicas y
% diversas bases de datos científicas como: IEEE SPRINGER LINK, GOOGLE SCHOLAR y
% LA UNIVERSIDAD TÉCNICA DE AMBATO, se ha recopilado una extensa variedad de
% información. Asimismo se limitó la búsqueda al período comprendido entre 2018 y
% 2023.

% \bigbreak
% Sin embargo, para asegurar la relevancia de los datos recopilados se
% aplicaron criterios de inclusión y exclusión detallados en la Tabla
% \ref{tabla_1}, descartando aquellos artículos que no cumplían con estos
% criterios.

% \begin{table}[h]
    \centering
    \caption{Criterios de inclusión y exclusión}
    \label{tabla_1}
    \begin{tabular}{|p{6cm}|p{6cm}|}
    \hline
    \textbf{Inclusión} & \textbf{Exclusión} \\
    \hline
    Artículos escritos en inglés y español. & Artículos escritos en otros idiomas que no sean en inglés y español. \\
    Artículos publicados entre los años 2018 – 2023. & Artículos publicados antes de 2018. \\
    Tesis y Artículos & Libros, páginas web y foros \\
    Artículo relacionados con monitoreo y control de riego usando Lógica Difusa. & Artículos no relevantes con la temática principal. \\
    Artículos disponibles en bibliotecas reconocidas o académicas que se enumeran en la sección de Base de datos. & Artículos disponibles en otras bibliotecas que no se enumeran en la sección de Base de datos. \\
    \hline
    \end{tabular}
    \end{table}
    

% Después de aplicar el proceso de exclusión, se identificaron un total de 21
% artículos relacionados con sistemas de riego que involucran la lógica difusa.
% Aunque estos artículos no coinciden directamente con el objetivo específico de
% esta investigación, se resaltan a continuación los trabajos más relevantes para
% el desarrollo del proyecto.


Para el desarrollo de esta investigación, es necesario llevar a cabo una revisión de trabajos previos vinculados al problema. Esta revisión tiene como objetivo identificar similitudes y diferencias con el proyecto actual. Como resultado de esta búsqueda, se han recopilado los siguientes estudios previos:

\bigbreak
En \cite{lema_holguin_implementacion_2018} se implementó un sistema automatizado de riego por goteo. Se utilizó una configuración de sensores para monitorear la humedad del sustrato, la conductividad eléctrica del suelo y el pH del agua, ubicados en contenedores de cultivo individuales con 24 plantas de tomate en cada uno. El proyecto presentó una solución para controlar y optimizar las condiciones de cultivo. Se implementó un invernadero que proporcionaba microclimas controlados para mitigar los impactos climáticos externos y mantener un entorno óptimo para el desarrollo de las plantas. Los sensores de temperatura, humedad, pH y conductividad eléctrica del sustrato se conectaron a dos controladores difusos tipo Mamdani, encargados de regular el riego y la administración de nutrientes en base a las lecturas obtenidas. En cuanto al hardware utilizado, se combinaron varios componentes específicos. Entre ellos se encuentran sensores como el Sensor Con-BTA y YL-68 para mediciones de pH y humedad, el sensor digital de temperatura DS18B20, junto con tarjetas Arduino Mega y Leonardo con Ethernet Shield para control y comunicación de datos. Se integraron bombas de agua, electroválvulas y otros sensores, cada uno asignado a funciones específicas dentro del sistema de riego por goteo automatizado. El diseño del sistema incluyó la definición de reglas difusas para los controladores, los cuales ejecutaban acciones en función de las lecturas de los sensores. Se realizaron mediciones periódicas mediante sensores y dispositivos de monitoreo para mantener las condiciones óptimas de cultivo. Los tiempos de respuesta de los controladores se mantuvieron dentro de los 5 segundos, comparados con tiempos obtenidos en Matlab. Sin embargo la presentación de datos numéricos que demuestren la estabilidad y precisión del sistema bajo diversas condiciones ambientales brindaría una comprensión más detallada de su rendimiento práctico.

\bigbreak
En \cite{alcivar_dominguez_sistema_2018} se implementó un sistema automatizado de riego basado en sensorización en los cultivos de ciclo corto. Este se compone de una variedad de componentes electrónicos clave. Entre ellos se encuentran sensores como el Módulo HL 69, el Módulo YL-83 y el DHT22, diseñados para recopilar datos precisos sobre la humedad del suelo, humedad ambiental y temperatura. El Arduino Mega se utilizó como el núcleo de control, encargado de coordinar la información de los sensores y tomando decisiones para el riego eficiente de los cultivos. Por otro lado, el software desarrollado para este sistema se basa en la plataforma de Visual Studio Community 2015. Esta elección el autor respalda por la versatilidad que ofrece, especialmente a través de Windows Form que facilitó el diseño y la presentación de consultas y reportes generados en el sistema. La integración de una base de datos MySQL-2012 le permitió la creación de archivos de tipo app config para manejar los parámetros de conexión al servidor, facilitando la lectura y almacenamiento de información para el sistema. La aplicación de escritorio se llevó a cabo bajo la metodología Cascada asegurando una secuencia metódica desde el análisis de necesidades hasta la implementación física y la validación del sistema en el terreno. Los resultados en la reducción del consumo de agua fueron de un 37.27\%, un aumento del 7.53\% en la productividad de los cultivos y un eficiente control de humedad del suelo. Si bien la elección de la metodología cascada para el desarrollo del software brindó una estructura metódica y secuencial en el proyecto de implementación del sistema de riego, se sugiere que podría haberse complementado con metodologías ágiles como Scrum o Kanban.

\bigbreak
En \cite{salazar_irrigation_2013} desarrolló un sistema de riego a través de Agentes Inteligentes usando la tecnología Arduino. La arquitectura que los autores propusieron se basa en agentes inteligentes que interactúan entre sí y con el entorno, utilizando una combinación de lógica difusa y el modelo BDI (Believe, Desire, Intention). El proceso inicia con la placa Arduino MEGA, seleccionada por el autor por su capacidad de conectividad y control, en comunicación con diversos sensores como el sensor de humedad del suelo, y actuadores que controlan el riego. Los sensores recopilan datos sobre la humedad del suelo, que son procesados por un Agente de Campo y luego enviados al Agente de Control. La lógica difusa la implementaron para utilizar un sistema MISO (entrada múltiple, salida única), datos de la humedad del suelo y su traducción en variables lingüísticas difusas. Esta arquitectura multiagente se desarrolla sobre la plataforma Java Agent DEvelopment Framework (JADE), siguiendo el estándar Foundation for Intelligent Physical Agents (FIPA) para sistemas multiagente. La interacción entre Arduino y la computadora se realiza mediante una conexión serial, y se ha desarrollado una interfaz de software que muestra detalles de los sensores y el estado del sistema de riego, con modos manual y automático. Las pruebas del sistema se llevaron a cabo en un proyecto de guayabas taiwanesas, demostrando que el sistema de riego basado en agentes inteligentes permitió un monitoreo en tiempo real de las necesidades de agua de cada cultivo, adaptándose a las condiciones específicas del suelo. Sin embargo, la omisión de detalles específicos sobre el tipo de hardware utilizado más allá de la placa Arduino MEGA limita la comprensión completa del sistema.

\bigbreak
En \cite{hasan_implementation_2018} los autores se centran en la gestión óptima de la humedad del suelo. Este sistema automatizado utiliza sensores especializados para monitorear constantemente el contenido de humedad en el suelo, enviando datos al sistema cada segundo. El sistema emplea lógica difusa para calcular el nivel de control requerido. Esto activa un método de control que ajusta una válvula en la línea de suministro de agua, permitiendo el flujo desde un tanque elevado estático hacia el campo. El proceso de riego se adaptaba dinámicamente según los niveles de humedad hasta al alcanzar el nivel óptimo. Este modelo de riego automatizado basado en un algoritmo difuso se encargaba de clasificar la humedad en seis clases: 'Muy bajo, bajo, medio, preciso, superior y desbordado'. Además, utiliza microcontroladores, motores y bombas para su funcionamiento. Los resultados que obtuvieron los autores de las pruebas comparativas entre el método de riego tradicional y este modelo propuesto fue una reducción del 12.3\% en el uso de agua. De esta manera sería favorable si se incluyeran detalles técnicos más específicos sobre los microcontroladores, motores y bombas utilizados, así como información detallada sobre la metodología de las pruebas comparativas entre el método tradicional de riego y el modelo propuesto.

\bigbreak
En \cite{munir_intelligent_2019}, desarrollaron un sistema de riego inteligente que se apoya en la lógica difusa y la tecnología blockchain para gestionar eficientemente el riego de las plantas. Utilizando datos provenientes de sensores que monitorean variables como temperatura, humedad, intensidad de luz, nivel de humedad del suelo y estado de la bomba de agua, este sistema toma decisiones en tiempo real respecto al riego de las plantas. Este sistema se compone de hardware como Arduino-UNO R3, Bread Board, sensores DHT-11 de temperatura y humedad, y sensor YL-69 de humedad del suelo. Estos dispositivos recopilan información crucial sobre las condiciones ambientales y del suelo, vital para determinar las necesidades hídricas de las plantas. La conectividad inalámbrica se logra mediante un módulo Wi-Fi ESP 8266-01, facilitando la transmisión de datos entre sensores y el servidor central. La aplicación móvil para Android permitió a los usuarios monitorear y controlar remotamente el sistema de riego, configurar horarios de riego, verificar la humedad del suelo y recibir recomendaciones personalizadas para el cuidado de las plantas. La metodología empleada se basa en Fuzzy Logic y Blockchain. La lógica difusa facilita la toma de decisiones sobre las necesidades de riego, mientras que la tecnología blockchain asegura la conectividad segura de dispositivos en un entorno IoT, permitiendo el acceso exclusivo a dispositivos confiables. El resultado que obtuvo el autor de las pruebas experimentales, la mayoría alcanzaron un 100\% de efectividad, es decir que la precisión general fue del 95.83\%. Sin embargo podría haberse reforzarse estos resultados con una mayor descripción de cómo se realizaron estas pruebas, qué variables se consideraron y cómo se evaluó la efectividad.

\bigbreak
En \cite{al-ali_iot-solar_2019}, se menciona un sistema de riego inteligente alimentado por energía solar para la agricultura. El sistema recopila múltiples tipos de datos incluyendo lecturas de humedad y temperatura del suelo, provenientes de sistemas de riego y componentes electrónicos. Los objetivos principales señalados por el autor se centran en el monitoreo remoto, control de humedad, temperatura, energía solar y riego automático en entornos agrícolas. Para poder lograr esto, desarrollaron algoritmos de control y lógica difusa, los cuales les permitió medir la eficiencia del riego, la adaptabilidad a condiciones climáticas variables y el consumo de energía solar. El sistema se basó en un controlador, myRIO, un procesador ARM de doble núcleo y una matriz FPGA, conectado a una variedad de sensores como humedad del suelo, humedad y temperatura, sensor de flujo, interruptores de flotador y actuadores como bombas de achique y de diafragma. El proceso de control se desglosa en tres modos: control local, monitoreo y control móvil, y control basado en lógica difusa. El último modo implica tomar decisiones sobre el encendido o apagado de las bombas basándose en lecturas de sensores y un algoritmo de control difuso. El sistema de riego inteligente propuesto se diseñó con la premisa de ser accesible y controlable desde cualquier ubicación y en cualquier momento, ofreciendo al usuario la capacidad de asumir el control total del sistema cuando así lo desee. Para lograr esta funcionalidad utilizaron un servidor web remoto que se respalda con una base de datos en Lenguaje de Consulta Estructurado (SQL). Esta base de datos se encarga de almacenar de manera organizada y estructurada las mediciones provenientes de los sensores utilizados en el sistema. Ademas, se utilizó un controlador de sistema en un chip de placa única con conectividad WiFi y conexiones a una celda solar para leer sensores de campo y emitir señales de comando para operar las bombas de riego. No obstante, sería útil incorporar detalles de resultados numéricos o métricas específicas como la reducción del consumo de agua o el porcentaje de mejora en la eficiencia del riego.

\bigbreak
En \cite{orozco_jaramillo_diseno_2019} se ha diseñado e implementado una red de sensores para monitorear los niveles de radiación solar en la ciudad de Loja, se destaca la utilización de sensores GUVA-S12SD, Raspberry Pi 3B y una aplicación móvil desarrollada en Android Studio para llevar a cabo la medición y visualización de estos niveles. El proceso de adquisición de datos se realiza a través de sensores especializados, como el sensor GUVA-S12SD, dedicados a monitorear los niveles de radiación solar. Estos sensores capturan la información y la envían a una unidad central, la Raspberry Pi, que actúa como concentrador de datos. Posteriormente, la Raspberry Pi recibe, procesa y almacena los datos provenientes de los sensores en una base de datos MySQL (SQL Structured Query Language). La conexión entre la base de datos MySQL y la aplicación móvil desarrollada en Android Studio permitió el acceso a los datos almacenados en tiempo real. La aplicación móvil disponía de secciones específicas para presentar esta información, incluyendo gráficos o mapas que representan los niveles de radiación solar recopilados, tanto en tiempo real como los datos históricos. Además, permitía a los usuarios realizar consultas a la base de datos para acceder a detalles históricos o información detallada sobre los niveles de radiación solar. El proceso de adquisición y almacenamiento de datos implica la conversión de datos binarios a valores hexadecimales y luego a valores decimales antes de su almacenamiento en la base de datos. Los autores llevaron a cabo pruebas de adquisición de datos en diferentes nodos sensores a lo largo de varios días y en diversas condiciones climáticas para validar la efectividad y precisión del sistema. De este modo podría haber señalado ha detalle los procesos de calibración o validación de los sensores utilizados para garantizar la precisión de las mediciones.

\bigbreak
En \cite{castillo_herrero_desarrollo_2020} se centró en el desarrollo de una aplicación móvil que utilizó tecnología NFC (Near Field Communications) para gestionar datos de sensores de temperatura asociados a una base de datos en un servidor cloud de IoT. El objetivo principal que señalaron los autores es permitir a los usuarios visualizar y acceder a la información de temperatura de manera rápida y segura. El proceso comienza con la captación de información a través de sensores de temperatura asociados a etiquetas NFC. Estos datos se almacenan y gestionan en un servidor dividido en dos secciones: una base de datos y un API REST (Application Programming Interface - Representational State Transfer) desarrollado con Python y Flask. La aplicación móvil desarrollada en Android Studio se comunica con este servidor para acceder a la información almacenada y presentarla a los usuarios de manera intuitiva. El acceso a la información varía según los roles de los usuarios. La lectura y escritura de datos en las etiquetas NFC permiten identificar sensores, registrar temperaturas y acceder al historial de registros. El sistema garantiza la seguridad mediante el uso de contraseñas y roles, asegurando que las acciones de modificación en la Base de Datos se realicen únicamente por usuarios autorizados. Además, se implementa una estrategia de actualización periódica del proyecto en Google Drive para preservar la integridad y continuidad del desarrollo de la aplicación. Sin embargo podría haber sugerido la inclusión de detalles adicionales sobre la eficacia y precisión de la tecnología NFC utilizada para la captación de datos de los sensores.

\bigbreak
En \cite{krishnan_fuzzy_2020} implementaron un sistema de riego inteligente basado en lógica difusa que utiliza Internet de las cosas, donde utilizaron la integración de varios sensores, incluyendo el sensor DHT11 para medir humedad y temperatura. Además utilizaron otros sensores para monitorear diversos parámetros del campo agrícola en tiempo real. Estos datos se muestran en una pantalla LCD (Liquid crystal display) y se transmiten al usuario a través de tecnología GSM (Global System For Mobile Communication) para facilitar el monitoreo remoto. La infraestructura hardware empleada en este sistema comprende un conjunto de sensores, como el sensor de humedad del suelo, sensor de temperatura, sensor de humedad, sensor de lluvia e incluso un sensor de imagen de hoja de planta. Todos estos sensores están conectados a un controlador Arduino, que actúa como nodo de dispositivo final. Un aspecto destacado es el uso de paneles solares para alimentar el sistema durante la disponibilidad de luz solar. La tecnología GSM se empleó para la automatización dando como resultado la disminución de la necesidad de trabajo manual. Los autores realizaron pruebas al sistema, donde se evaluó de manera comparativa los métodos tradicionales de riego, el riego por goteo e inundación manual. En una comparación con el riego por goteo y el riego manual por inundación, se destacó que el sistema de riego inteligente propuesto bombea agua en un período de 7 horas, mientras que los métodos tradicionales lo hacen en 12 y 20 horas respectivamente. Es decir que el sistema de riego inteligente utiliza el motor durante solo el 9,72\% del tiempo total de riego, en comparación con el 16,67\% y 27,78\% utilizados por el riego por goteo y el riego manual por inundación respectivamente. En mi opinión podrían haberse centrado en profundizar en la comparación de la eficiencia del sistema propuesto con los métodos tradicionales de riego, proporcionando más detalles sobre la precisión de las mediciones de humedad y temperatura realizadas por los sensores.

\bigbreak
En \cite{kumar_irrigation_2020} abordo el desarrollo de un sistema de control de riego basado en la utilización de cuatro tipos de sensores: temperatura, intensidad de luz, humedad del suelo y humedad. Estos sensores fueron necesarios para en el monitoreo y control del proceso de riego en entornos agrícolas. El objetivo principal que presentaron los autores de este proyecto es mejorar la eficacia del riego a través de un enfoque que emplea el sistema ANFIS (Sistemas Adaptativos de Inferencia Neuro-Difusa) en una Red de Sensores Inalámbricos (WSN). Se implementó el protocolo PEGASIS para la recolección eficiente de datos en el sistema de riego, estableciendo una cadena entre los nodos para transmitir la información a la base. El enfoque ANFIS se utilizó para la toma de decisiones en el riego, activando automáticamente la bomba de agua y la lámpara según las condiciones detectadas por los sensores. Posteriormente se inició con la recolección de datos, asegurando una conexión continua y una recolección eficiente. ANFIS integró 81 reglas difusas para el sistema de riego automático, de manera que mapeaba estas variables para la activación de la bomba de agua y la lámpara según las condiciones detectadas por los sensores. En términos de pruebas y resultados los autores compararon el método ANFIS-PEGASIS con otras técnicas existentes. Se evaluaron parámetros como el rendimiento del retardo de extremo a extremo, consumo de energía y rendimiento del sistema. Asimismo los autores mencionan que los resultados de la simulación demostraron que el método propuesto (ANFIS-PEGASIS) superó a otras técnicas existentes en términos de E2ED, rendimiento y sostenibilidad agrícola. Sería muy útil para el estudio si se incluyera la comparación detallada que se realizó con otras técnicas.

\bigbreak
En \cite{mohammed_intelligent_nodate} presentó un sistema de riego inteligente diseñado para la región oriental de Marruecos, en el cual se llevó a cabo la recopilación de datos mediante sensores estratégicamente ubicados en una región agrícola. Los datos se adquirieron en tiempo real a través de una estación meteorológica y sensores de humedad del suelo conectados a las plantas. Se implementó un controlador lógico difuso basado en reglas Mamdani que evaluaba la humedad del suelo, temperatura y radiación solar para determinar el tiempo óptimo de riego. El sistema se complementó con una placa Arduino Uno R3, permitiendo la adquisición de datos mediante sensores analógicos y digitales. Se realizaron simulaciones con datos climáticos reales para validar el sistema y analizar su comportamiento en diferentes estaciones del año. Se seleccionó un manzano como planta de prueba, y los resultados mostraron que el sistema mantuvo la humedad del suelo por encima del umbral deseado, evitando el riesgo de riego insuficiente. A pesar de ello, sería útil profundizar en la descripción técnica del hardware utilizado para mejorar la comprensión del sistema.

\bigbreak
En \cite{benyezza_zoning_2021} se centró en maximizar la eficiencia hídrica y energética en la agricultura. Se implementó un sistema de riego basado en tecnología de control difuso e IoT, el cual utilizó sensores de humedad del suelo y temperatura ambiental en zonas estratégicas del invernadero. Estos sensores formaron una red inalámbrica que transmitió datos a un servidor Node-RED para su procesamiento. El proceso se fundamentó en la tecnología de control difuso para tomar decisiones óptimas de riego, sin requerir un modelo matemático preciso del suelo. Se implementó un Controlador Lógico Difuso (FLC) en un nodo central, basándose en reglas difusas extraídas de conocimientos expertos y experiencias previas. El hardware empleado incluyó Arduino Nano, módulo nRF24L01þ para comunicación inalámbrica, sensor DHT22 para medir la temperatura y una batería para la alimentación de los nodos. Además, se aplicaron medidas de seguridad en el editor Node-RED para asegurar el sistema contra accesos no autorizados. En las pruebas comparativas con otros métodos de riego, los autores observaron que el sistema propuesto logró un consumo reducido de agua, significativos ahorros en el costo de producción, una disminución en el consumo energético y un control efectivo de la humedad del suelo, favoreciendo el crecimiento de las plantas. Se destacó que este sistema inteligente de riego, al incorporar la irrigación por zonas, control difuso, comunicación inalámbrica y control remoto, ofreció mejoras superiores al 80\% en comparación con estrategias de riego tradicionales. Sin embargo, podría haber estado en evaluación y adaptaciones continuas del sistema para abordar una gama más amplia de condiciones ambientales. Al ampliar la consideración de diversos escenarios, podría haberse optimizado aún más el rendimiento del sistema, aumentando su versatilidad y eficacia en una variedad más amplia de entornos agrícolas.

\bigbreak
En \cite{widyawati_fuzzy_2022}, se desarrolló un diseño de lógica difusa para controlar la duración del tiempo de riego en un invernadero. Este diseño empleo datos obtenidos a partir de sensores de humedad del suelo y el sensor DHT22 para monitorizar la humedad y temperatura. El objetivo central se centró en el diseño y la implementación de un Controlador de Lógica Difusa (FLC) tipo Mamdani, integrado en una placa NodeMCU ESP8266. Este sistema utilizó los datos de los sensores como entrada para determinar la duración óptima del riego. Las pruebas y evaluaciones comparaban la lógica difusa en NodeMCU ESP8266 con el software MATLAB (Matrix Laboratory) para validar su desempeño. La lógica difusa se estructuraba en etapas de fuzzificación, inferencia basada en reglas y defuzzificación para determinar la duración del riego a partir de la humedad y temperatura del suelo. Los valores de entrada se convirtieron en datos lingüísticos difusos, aplicando funciones de pertenencia difusas y reglas de conexión para determinar la duración óptima del riego. La implementación del sistema integró una puerta de enlace y nodos de sensores, donde la lectura y procesamiento de datos se realizaron mediante lógica difusa en NodeMCU ESP8266, programada en lenguaje C con Arduino IDE. El artículo detalla los componentes de hardware empleados, entre los que se encuentran NodeMCU ESP8266, sensores de humedad del suelo, sensor DHT22, módulo TP4056, batería de ion de litio 18650, Raspberry Pi y cables jumper. Los resultados de las pruebas comparativas entre la lógica difusa basada en NodeMCU ESP8266 y la de MATLAB indican un error promedio del 0.59\%. Esto sugirió que el controlador de lógica difusa con el método Mamdani integrado en el NodeMCU ESP8266 ofrece una alta precisión para controlar la duración del riego en invernaderos. Sin embargo, se podría considerar la inclusión de análisis y resultados adicionales en términos de la robustez del sistema frente a variaciones extremas en las condiciones ambientales. Esto podría proporcionar una visión más completa de la aplicabilidad del sistema.

\bigbreak
En \cite{haiyan_intelligent_2022} se destacó una plataforma avanzada de gestión de recursos hídricos para agricultura que integra datos de temperatura, humedad y crecimiento de plantas mediante sensores y cámaras digitales. Esta plataforma emplea técnicas de procesamiento de imágenes digitales y lógica difusa para analizar la información de crecimiento de cultivos, evaluar la necesidad de riego y detectar la escasez de agua en las plantas. El proceso de adquisición de datos inicio con la detección precisa del suelo y su entorno mediante sensores especializados, como el sensor de humedad del suelo SM3002B y el sensor de temperatura y humedad del aire SHT30. Estos dispositivos capturaron información en tiempo real sobre las condiciones del suelo y el ambiente de cultivo. Posteriormente, esta información se transmite de forma inalámbrica a un servidor central utilizando un protocolo específico para su procesamiento. El procesamiento central se lleva a cabo mediante microcontroladores como la microcomputadora de un solo chip modelo STC15F2K60S2 y el chip MSP430. Estos componentes realizan la conversión y procesamiento de los datos recolectados. Además, la plataforma cuenta con una pantalla LCD de 320 * 240, que muestra información ambiental en tiempo real. Este sistema emplea algoritmos de control difuso y control PID (Control Proporcional Integral Derivativo) para ajustar automáticamente el riego según la información recopilada. La alimentación de la plataforma se sustenta en energía solar. En relación con la estabilidad y confiabilidad de la red de comunicación, los autores realizaron un experimento que se enfoca en comparar las redes de transmisión punto a punto y las redes de múltiples saltos. Este estudio destacó la importancia de una adecuada distribución de nodos y la gestión de la distancia entre ellos para mantener una comunicación eficaz y estable. Estos resultados son fundamentales en sistemas de riego, donde la fiabilidad de la comunicación es esencial para el control efectiva del sistema.

\bigbreak
En \cite{noauthor_fuzzy_2023} presentó un sistema de control de riego basado en lógica difusa, destinado a optimizar la producción agrícola mediante el uso de datos meteorológicos y datos de sensores especializados, como sensores de humedad del suelo, temperatura y humedad. Este sistema se integra con sensores en robots móviles, datos meteorológicos y una aplicación móvil, los cuales monitorean constantemente variables clave como la humedad del suelo, la humedad y la temperatura. El proceso se inició con la recopilación y envío de datos de estos sensores a un servidor en la nube. Estos datos se almacenaron en una base de conocimientos previamente establecida sobre las necesidades específicas de distintos cultivos. Los autores utilizaron reglas difusas para la toma decisiones sobre el funcionamiento del motor del sistema de riego, basándose en las condiciones medidas por los sensores. El sistema se apoya en sensores como el DTH11 para la medición de temperatura y humedad, así como sensores de humedad del suelo. Además, utiliza hardware como Arduino Mega, controlador ICL2930 y microcontrolador ESP8266, capaz de comunicarse por WiFi (Wireless Fidelity), para el procesamiento y la comunicación de datos. La aplicación móvil proporciono visualización de parámetros del sensor y control remoto del sistema de riego, permitiéndole al agricultor supervisar y en algunos casos influir en el proceso de toma de decisiones. Las pruebas que realizaron los autores para evaluar el sistema incluyeron diversos aspectos. Desde la evaluación de las reglas difusas hasta la precisión en la toma de decisiones para activar o desactivar el motor de riego en varios escenarios. Se menciono una precisión del 97\% en la toma de decisiones durante la evaluación, lo que significa que, de 1079 casos de prueba, el sistema de inferencia difusa evaluó correctamente la salida para 1046 muestras. Sin embargo, seria de mayor utilidad la profundización en la discusión sobre la precisión del sistema en condiciones extremas.

\bigbreak
El estudio presentado en \cite{neugebauer_fuzzy_2023} se centró en el análisis de un sistema de riego inteligente basado en lógica difusa. Este análisis combinó datos teóricos, mediciones reales y simulaciones computacionales. Utilizaron el Método de Elementos Finitos (FEM) para modelar la propagación del agua en el suelo, permitiendo simular la irrigación. Se evalúo múltiples variables como la temperatura ambiente, la humedad del suelo y la hora del día, las cuales controlan la intensidad del riego a través de FLC. Mediante un modelo 2D, lograron simular y comparar distintos sistemas de riego, facilitando la elección del más eficiente para suministrar agua a las raíces de las plantas con el menor consumo posible. El uso de un controlador de lógica difusa, basado en variables como temperatura, humedad del suelo y hora del día, permite determinar el tiempo y la cantidad óptima de riego. Las pruebas incluyeron la ubicación de sensores de humedad del suelo a diferentes profundidades y su impacto en el consumo de agua. Además, se realizaron pruebas comparativas entre sistemas de riego convencionales y aquellos controlados por lógica difusa bajo diversas condiciones y distribuciones de variables. En estas pruebas se observó que el sistema controlado por lógica difusa utilizó un 13\% menos de agua en comparación con el sistema convencional. Sin embargo, en una simulación con distribuciones diferentes de variables pero con las mismas reglas, se notó un aumento considerable en el consumo de agua por parte del sistema controlado por lógica difusa. Sin embargo, sería valioso proporcionar resultados más detallados obtenidos de las comparaciones entre diferentes sistemas de riego controlados por lógica difusa.

\bigbreak
En \cite{ramos_galindo_diseno_2023} se basó en el desarrollo de una aplicación móvil utilizando el framework Flutter para el monitoreo de humedad del suelo. Este sistema se complementa con un conjunto de sensores controlados por un microcontrolador ESP32 y hace uso de Firebase como base de datos en tiempo real. La arquitectura utilizada implica el uso de Flutter como framework principal para la creación de la aplicación móvil, permitiendo el desarrollo de interfaces intuitivas y la visualización de datos en tiempo real provenientes de sensores. El microcontrolador ESP32 se elige por su capacidad integral para manejar las necesidades del proyecto, incluyendo la adquisición de datos analógicos de los sensores de humedad. En cuanto a la base de datos y autenticación, Firebase es seleccionado por su capacidad de mantener operativa una base de datos en tiempo real las 24/7. Esto permitió almacenar y gestionar los datos obtenidos por los sensores de manera eficiente y accesible. El proceso inicio con la adquisición de datos a través de sensores de humedad conectados al microcontrolador ESP32, que capturan información en tiempo real sobre las condiciones del suelo. Estos datos se transmiten a la base de datos en Firebase para su procesamiento y almacenamiento. La aplicación móvil desarrollada en Flutter permite a los usuarios visualizar esta información de manera intuitiva y tomar decisiones basadas en las necesidades específicas de cada cultivo. La arquitectura general del sistema comprende el uso de sensores conectados al microcontrolador ESP32 para la adquisición de datos, Firebase para el almacenamiento en la nube, y Flutter para la interfaz de usuario y visualización de datos en tiempo real. Este sistema busca proporcionar una solución integral para el monitoreo de humedad del suelo, permitiendo a agricultores, horticultores y entusiastas del cultivo acceder fácilmente a información valiosa para mejorar el cuidado de sus cultivos. Sin embargo, podría haber brindado una mayor exploración en los detalles técnicos de la integración de los sensores con el microcontrolador ESP32 y la transmisión de datos a Firebase.

\bigbreak
Tras revisar diversos antecedentes, se destaca cómo la lógica difusa ha tenido un gran impacto al interpretar datos complejos, como la humedad, la temperatura y los niveles de humedad. Esta técnica resulta especialmente útil para representar estados en motores y sistemas difusos, sobre todo en situaciones donde los datos de entrada y salida no son precisos. Además, se ha observado un avance considerable en la precisión y diversidad de sensores, desde aquellos diseñados para medir la humedad del suelo, como el Módulo HL 69, SM3002B y el sensor YL-69, hasta sensores específicos para la lluvia, como el sensor de lluvia y el sensor de gota de lluvia YL-83.
\bigbreak
A pesar de la relevancia que muestran los estudios revisados sobre la aplicación de la lógica difusa en sistemas y sensores, su enfoque no se centra en su implementación directa en el desarrollo de aplicaciones móviles específicamente para el control y monitoreo de riego. Es importante destacar que a pesar de la importancia que tiene el riego de agua en viveros para el cuidado de las plantas, los estudios revisados no abordan específicamente esta práctica en ese entorno. Esta falta de información resalta la necesidad de profundizar en la implementación de sistemas de riego en viveros.

\section{Fundamentación teórica}

Esta sección se basa en una exhaustiva revisión de antecedentes, con el
propósito de proporcionar los fundamentos y teorías necesarios para respaldar y
contextualizar la investigación. Su objetivo es establecer un marco sólido que
sustente cada fase del estudio.

\subsection{Vivero}
Un vivero es un espacio destinado a la producción y cuidado de plantas, especialmente para su venta \cite{peris_para_nodate}. En estos lugares se maneja una amplia gama de plantas, incluyendo ornamentales, forestales y frutales. Al revisar investigaciones previas, se nota que se mencionan sistemas de riego en entornos como invernaderos, granjas y otros lugares, pero no específicamente en viveros. Sin embargo, el riego es esencial en los viveros para garantizar el crecimiento saludable de las plantas. Dado que los antecedentes no detallan la implementación de sistemas de riego en viveros, es crucial resaltar esta omisión.

\subsection{Lógica difusa}
La lógica difusa es una técnica que se utiliza para la toma de decisiones con datos incompletos. Esta técnica se basa en el uso de grados de verdad, en lugar de la dicotomía verdadero o falso de la lógica clásica. En \cite{lema_holguin_implementacion_2018} \cite{salazar_irrigation_2013} \cite{hasan_implementation_2018} \cite{munir_intelligent_2019} \cite{al-ali_iot-solar_2019} \cite{krishnan_fuzzy_2020} \cite{kumar_irrigation_2020} \cite{benyezza_zoning_2021} \cite{noauthor_fuzzy_2023} \cite{neugebauer_fuzzy_2023} y \cite{mohammed_intelligent_nodate}  se utilizó para calcular parámetros de entrada, como la humedad del suelo, la temperatura y los niveles de humedad, generando resultados que representaban el estado del motor.

\subsection{Sistemas de control}
Son conjuntos de dispositivos mecánicos y/o electrónicos que regulan, administran y controlan otros dispositivos o sistemas mediante lazos de control. Estos sistemas se utilizan para proporcionar la respuesta deseada controlando la salida de un proceso. La lógica difusa tipo Mamdani se utilizo en \cite{mohammed_intelligent_nodate} \cite{widyawati_fuzzy_2022} para modelar sistemas difusos y tomar decisiones basadas en reglas lingüísticas, lo que la hace adecuada para sistemas en los que las entradas y salidas no son precisas. Por otro lado, el algoritmo PID utilizado en \cite{haiyan_intelligent_2022} es un método de control de retroalimentación clásico que se utiliza para corregir errores o fluctuaciones en un sistema. Ambos tipos de algoritmos se utilizan en sistemas de control para mejorar el rendimiento y la precisión en el control de sistemas dinámicos.

\subsection{Sensores}
Los sensores representan dispositivos fundamentales en diversos campos al detectar o medir alteraciones en su entorno y transformar dicha información en señales eléctricas o datos.

\subsection*{Sensores de Humedad y Temperatura}
Los sensores de humedad y temperatura, como el DTH11, el DHT22 y el DS18B20 utilizados en \cite{lema_holguin_implementacion_2018} \cite{alcivar_dominguez_sistema_2018} \cite{munir_intelligent_2019} \cite{krishnan_fuzzy_2020} \cite{benyezza_zoning_2021} \cite{widyawati_fuzzy_2022} y \cite{noauthor_fuzzy_2023}, desempeñan el papel necesario para medir simultáneamente la humedad relativa y la temperatura en un entorno. Cada uno de estos sensores tiene sus propias características distintivas. Mientras que el DTH11 destaca por su costo económico y precisión aceptable, el DHT22 ofrece una mayor exactitud en las mediciones y un rango de operación más amplio. Por otro lado, el DS18B20 se destaca por su interfaz digital y su precisión en mediciones de temperatura. En \cite{castillo_herrero_desarrollo_2020}, los sensores de humedad fueron elementos necesarios para obtener el acceso a los datos mediante la tecnología NFC.

\subsection*{Sensores de Humedad del suelo}
Los sensores de humedad del suelo son dispositivos diseñados para medir el contenido de humedad en la tierra. Entre los tipos de sensores comunes se encuentran el Módulo HL 69, el SM3002B y el sensor YL-69 los cuales se usaron en \cite{alcivar_dominguez_sistema_2018} \cite{haiyan_intelligent_2022} y \cite{haiyan_intelligent_2022}. Cada uno posee particularidades distintivas que los diferencian en su aplicación. El Módulo HL 69 destaca por su versatilidad y capacidad para medir la humedad en diferentes tipos de suelo, ofreciendo una amplia gama de aplicaciones. Por otro lado, el SM3002B se caracteriza por su precisión en la medición y su resistencia a condiciones ambientales desfavorables. Por último, el sensor YL-69 es conocido por su sencillez y costo accesible, siendo una opción práctica para aplicaciones básicas de monitoreo de humedad en el suelo.

\subsection*{Sensores de Lluvia}
Un sensor de lluvia es un dispositivo diseñado para detectar la presencia o la intensidad de la lluvia. Entre los tipos comunes se encuentran el sensor de lluvia y el sensor de gota de lluvia YL-83. Estos sensores presentan características particulares que los diferencian. El sensor de lluvia referenciado en \cite{krishnan_fuzzy_2020}, se enfoca en utilizar métodos basados en lógica difusa para la detección y el análisis de la lluvia. Por otro lado, el sensor de gota de lluvia YL-83, mencionado en \cite{alcivar_dominguez_sistema_2018}, se destaca por su capacidad para detectar gotas individuales de lluvia y su sensibilidad para medir la intensidad de la misma.

\subsection{Microcontrolador}
Un microcontrolador es un chip que integra procesador, memoria y periféricos, utilizado para controlar operaciones en tiempo real. Entre ellos, el ESP8266 se destaca por su conectividad Wi-Fi el cual se uso en \cite{hasan_implementation_2018} y \cite{widyawati_fuzzy_2022}, el STC15F2K60S2 ofrece control preciso usado en \cite{haiyan_intelligent_2022}, y el ESP32 es más avanzado, no solo tiene potencia de procesamiento sino también soporte para Wi-Fi y Bluetooth. Además, en \cite{ramos_galindo_diseno_2023} el ESP32 se ha utilizado para recopilar datos de sensores, procesar la información y transferirla en tiempo real a una base de datos a través de Wi-Fi, aprovechando sus capacidades avanzadas para tareas simultáneas y conectividad inalámbrica.

\subsection{Controlador}
Un controlador es un dispositivo electrónico diseñado para regular y dirigir el funcionamiento de otros componentes en un sistema. En el caso del ICL2930 utilizado en \cite{hasan_implementation_2018}, se destaca por su capacidad para manejar motores, ofreciendo una interfaz eficiente y precisa para controlar la velocidad y dirección de motores DC.

\subsection*{Arduino}
Arduino es una plataforma de hardware de código abierto para proyectos electrónicos. Entre sus modelos, el Arduino Mega usado en \cite{lema_holguin_implementacion_2018} \cite{alcivar_dominguez_sistema_2018} y \cite{hasan_implementation_2018} se destaca por su mayor cantidad de pines, ideal para proyectos complejos. El Arduino-UNO R3 es versátil y ampliamente utilizado, perfecto para aplicaciones estándar, este fue utilizado en \cite{haiyan_intelligent_2022} y \cite{mohammed_intelligent_nodate} . Mientras que el Arduino Nano ofrece funcionalidad en un formato compacto, ideal para proyectos con limitaciones de espacio, este se utilizó en \cite{benyezza_zoning_2021}.

\subsection*{Raspberry Pi}
El Raspberry Pi es una serie de miniordenadores de placa única desarrollados para promover la enseñanza de informática y la creación de proyectos de bajo costo. Entre los modelos de Raspberry Pi, el Raspberry Pi 3 es destacado por su rendimiento superior con un procesador más rápido y mejoras en la conectividad. Este modelo ha sido utilizado en diversos contextos, como se cita en \cite{benyezza_zoning_2021} y \cite{jaramillo_diseno_2020}.

\subsection{Metodología Cascada}
La metodología en cascada es un enfoque secuencial y lineal para el desarrollo de software, donde cada fase del proyecto se completa antes de pasar a la siguiente. En \cite{alcivar_dominguez_sistema_2018}, se utilizó para asegurar un proceso de desarrollo sistemático y ordenado. Esta metodología permite una planificación detallada, la estimación precisa de costos y plazos, y establece una secuencia lógica en las etapas de creación del sistema de riego automatizado. Cada etapa se ejecuta de manera secuencial, comenzando desde la definición de requisitos, pasando por el diseño, la implementación, pruebas y finalmente la entrega. Este enfoque permite una estructura clara y una gestión eficiente de cada fase del proyecto.


\subsection{Metodología XP}
La metodología Extreme Programming (XP) es un enfoque ágil de desarrollo de software que se centra en la mejora continua, la adaptabilidad y la entrega rápida de software de alta calidad \cite{cabezas_orellana_alisis_2023}. XP se basa en una serie de prácticas y valores diseñados para maximizar la satisfacción del cliente y la eficiencia del equipo de desarrollo. Es relevante mencionar que, al revisar los antecedentes existentes, no se especifica la metodología utilizada en los proyectos previos. Dada la falta de información sobre las metodologías previas, se ha decidido adoptar la metodología XP para el desarrollo del presente proyecto.


\subsection{Base de Datos}
Una base de datos es un sistema organizado para almacenar y gestionar datos de manera estructurada. Entre sus tipos está la base de datos SQL, que utiliza el lenguaje SQL para administrar datos, como se evidencia en \cite{al-ali_iot-solar_2019}. MySQL es un sistema de gestión de bases de datos relacional ampliamente empleado en diversos sistemas, según lo señalado en \cite{alcivar_dominguez_sistema_2018} y \cite{castillo_herrero_desarrollo_2020}. Por último, Firebase que es una plataforma de Google, que proporciona servicios de base de datos en tiempo real, siendo utilizado para mantener una base de datos en tiempo real, como se menciona en \cite{ramos_galindo_diseno_2023}, permitiendo la sincronización entre dispositivos y aplicaciones.


\subsection{Tecnología NFC}
La tecnología NFC, conocida como Near Field Communication, se emplea por su habilidad para establecer conexiones entre dispositivos electrónicos de manera sencilla e intuitiva. En \cite{haiyan_intelligent_2022} cada sensor está asociado con una etiqueta NFC específica. Cuando un usuario requiere acceder al historial de temperaturas de un sensor en particular, simplemente necesita escanear la etiqueta NFC correspondiente usando esta tecnología. Esta forma de interacción simplifica el proceso de obtención de datos sensoriales, ya que la lectura de la etiqueta NFC relacionada con un sensor específico desencadena la recuperación de información directamente en la aplicación.


\subsection{Arquitectura}
La arquitectura, en el contexto del desarrollo de software, se refiere a la estructura general y el diseño de un sistema o una aplicación. El Modelo Vista Controlador (MVC) es un patrón arquitectónico que organiza la lógica de una aplicación en tres componentes claves. El modelo que gestiona los datos y la lógica. La vista que maneja la presentación y la interfaz de usuario. Y finalmente el controlador que actúa como intermediario y gestiona las interacciones entre el modelo y la vista \cite{bascon_pantoja_patron_2004}. Al revisar los antecedentes existentes, no se encontró mención alguna sobre la arquitectura utilizada en los proyectos previos. Dada esta falta de información, se ha optado por la adopción del patrón MVC para la estructuración del proyecto.


\subsection{Arduino IDE}
Arduino IDE (Integrated Development Environment) es un entorno de desarrollo integrado utilizado para programar placas de hardware Arduino. Es un software que proporciona todas las herramientas necesarias para escribir, compilar y cargar código en las placas Arduino. En \cite{widyawati_fuzzy_2022} se utilizó para programar en lenguaje C el nodo sensor.


\subsection{Matlab}
Es un software utilizado para cálculos numéricos, simulaciones y análisis de datos en campos como ingeniería y ciencias. En \cite{widyawati_fuzzy_2022} se empleó MATLAB para evaluar y ajustar el rendimiento de la lógica difusa implementada en el dispositivo NodeMCU ESP8266. Este software se utilizó en este contexto particular para realizar análisis y pruebas relacionadas con la lógica difusa en el mencionado dispositivo.


\subsection{Flask}
Python es un lenguaje de programación interpretado, multiparadigma y multiplataforma, que ha ganado gran popularidad gracias a su legibilidad y su versatilidad en una amplia gama de aplicaciones, incluyendo el desarrollo web. Flask se ah usado en \cite{noauthor_fuzzy_2023} ya que es un framework de aplicaciones web escrito en Python, que se destaca por su simplicidad y flexibilidad en la construcción de estas aplicaciones. La elección de Python en el desarrollo web se debe a su facilidad de uso, su amplia aceptación en la comunidad de desarrolladores y su capacidad para integrarse con otros lenguajes y herramientas. Además, la ligereza y facilidad de aprendizaje de Flask lo hacen ideal para un desarrollo ágil en el contexto de aplicaciones web.

\subsection{Visual Studio Community}
Es una versión gratuita y robusta del Integrated Development Environment (IDE) de Microsoft, diseñado para desarrolladores que trabajan en aplicaciones web, móviles y de escritorio. Ofrece un conjunto de herramientas completo que incluye editores de código, depuradores, compiladores y más, facilitando la creación y el desarrollo de software en varios lenguajes de programación como C\#, C++, Python, entre otros. En \cite{alcivar_dominguez_sistema_2018} permitió a los desarrolladores tener las herramientas necesarias para crear una aplicación consola. Esta aplicación fue especialmente diseñada para analizar los datos recolectados por el sistema automatizado de riego.

\subsection{Aplicación Móvil}
Las aplicaciones móviles son soluciones diseñadas para proporcionar acceso a información y servicios a través de dispositivos móviles como teléfonos inteligentes y tabletas, como se menciona en \cite{noauthor_fuzzy_2023}. Android Studio es el IDE líder para aplicaciones Android, ya que ofrece herramientas que simplifican la creación y optimización de aplicaciones. En \cite{castillo_herrero_desarrollo_2020} se destaca su posición predominante en el desarrollo de aplicaciones para el sistema operativo Android, convirtiéndose en la elección natural para proyectos de esta índole. Mientras que en Flutter es un marco de trabajo de código abierto desarrollado por Google. En \cite{ramos_galindo_diseno_2023} se optó por su capacidad para crear aplicaciones multiplataforma con un alto rendimiento, su facilidad de desarrollo, la capacidad de personalización de la interfaz de usuario y el respaldo de una sólida comunidad de desarrolladores.

% \subsection*{Bombas de Agua}
% En \cite{hasan_implementation_2018} se emplean para trasladar agua de un lugar a otro, lo que le permite incrementar la presión en el suministro de agua o para fines específicos, como el riego de cultivos. Su aplicación responde a la necesidad de transferir agua de manera eficiente y controlada, siendo esencial en diversas áreas, desde la agricultura hasta el abastecimiento de agua potable. El diseño y la funcionalidad de estas bombas les permiten realizar estas tareas de manera efectiva,

% \subsection*{Chip MSP430}
% En \cite{haiyan_intelligent_2022} se utilizó porque cuenta con una serie de componentes integrados, opera con un amplio rango de voltajes y consume poca energía. Además, su diseño y conjunto de instrucciones lo hacen ideal para sistemas de control en tiempo real y aplicaciones de dispositivos integrados. Es muy utilizado en el desarrollo de sistemas electrónicos debido a su compatibilidad con diferentes entornos de desarrollo. Su disponibilidad generalizada y el soporte que ofrece lo convierten en una opción atractiva para proyectos que requieren un microcontrolador confiable y fácil de programar.

\section{Objetivo general}
Implantar un prototipo del sistema de monitoreo y control basado en lógica difusa para la gestión del consumo de agua en los cultivos del Vivero Michita.

\section{Objetivos específicos}
\begin{itemize}
    \item Analizar los procesos de consumo de agua en el Vivero Michita para mejorar la gestión de este recurso.
    \item Investigar herramientas y algoritmos para sistemas de monitoreo y control basados en lógica difusa.
    \item Desarrollar un prototipo del sistema de monitoreo y control basado en lógica difusa para gestionar el consumo de agua en los cultivos del Vivero Michita.
\end{itemize}

\chapter{METODOLOGÍA}
\section{Materiales}

\section{Métodos}
\subsection{Modalidad de investigación}
En esta sección, se proporcionará una descripción detallada de las modalidades de investigación adoptadas para el proyecto. A continuación se detalla cada una.

\bigbreak
\textbf{Investigación cuantitativa}\\
Este proyecto se apoya en una modalidad de investigación cuantitativa al emplear variables numéricas para llevar a cabo procesos estadísticos y presentar datos. Su enfoque se centra en la medición precisa de variables como niveles de humedad, temperaturas y patrones de riego, con el propósito de establecer relaciones cuantificables que mejoren la gestión del agua.

\bigbreak
\textbf{Investigación experimental}\\
Se plantea una investigación de tipo experimental para este proyecto, ya que se busca evaluar la efectividad y la precisión del sistema de gestión de riego. Esto implica la realización de pruebas controladas en viveros, variando los niveles de riego de acuerdo con los datos obtenidos por la lógica difusa.

\bigbreak
\textbf{Investigación aplicada}\\
Se trata de una investigación aplicada, ya que se utilizarán los conocimientos adquiridos a lo largo de la carrera para el desarrollo y ejecución de la propuesta. Esto implica la aplicación práctica de teorías, métodos y herramientas aprendidas en un entorno real para la creación del sistema de gestión de riego en el prototipo de vivero.


\chapter{RESULTADOS Y DISCUSIÓN}

\chapter{CONCLUSIONES Y RECOMENDACIONES}


%C. Materiales de referencia
\bibliographystyle{biblio/IEEEtran}
\bibliography{biblio/main}
\addcontentsline{toc}{chapter}{\bfseries REFERENCIAS BIBLIOGRÁFICAS}

\chapter*{Anexos}
\addcontentsline{toc}{chapter}{\bfseries\uppercase{Anexos}}

% \addtocontents{toc}{\protect\setcounter{tocdepth}{0}}
\appendix

\renewcommand{\thetable}{\theappendixs\arabic{table}}
\renewcommand{\thefigure}{\theappendixs\arabic{figure}}
\renewcommand{\thelstlisting}{\theappendixs\arabic{lstlisting}}
\captionsetup{list=no}

% \appendixs{Guía de entrevista aplicada a los trabajadores del Vivero Michita}
% \label{app:guia-entrevista-vivero-michita}
% \begin{table}[h]
% 	\caption{Guía de entrevista aplicada a los trabajadores del Vivero Michita}\label{tab:entrevista-vivero-michita}
% 	\begin{center}
% 		\begin{tabular}[c]{p{0.3cm} p{5.5cm} p{3.5cm} p{3.5cm}}
% 			\hline
% 			\multicolumn{1}{c }{\textbf{N.}} & \multicolumn{1}{c}{\textbf{Pregunta}}                                                                                          & \multicolumn{1}{c}{\textbf{Respuesta}} & \multicolumn{1}{c}{\textbf{Observación}} \\
% 			\hline
% 			\addlinespace
% 			1                                & ¿Cuáles son los principales desafíos que enfrentan actualmente en cuanto al consumo de agua en los cultivos del vivero?        &                                        &                                          \\
% 			\addlinespace
% 			2                                & ¿Qué métodos o herramientas utilizan actualmente para monitorear y controlar el consumo de agua en los cultivos?               &                                        &                                          \\
% 			\addlinespace
% 			3                                & ¿Qué métodos de riego utilizan actualmente y con qué frecuencia?                                                               &                                        &                                          \\
% 			\addlinespace
% 			4                                & ¿Cuáles son los principales criterios que tienen en cuenta al decidir cuándo y cuánto regar los cultivos?                      &                                        &                                          \\
% 			\addlinespace
% 			5                                & ¿Cómo determinan las necesidades de agua de los diferentes tipos de plantas en el vivero?                                      &                                        &                                          \\
% 			\addlinespace
% 			6                                & ¿Qué problemas o desafíos han experimentado en el pasado con respecto al riego de los cultivos?                                &                                        &                                          \\
% 			\addlinespace
% 			7                                & ¿Existe algún tratamiento previo al agua antes de su uso en el riego? En caso afirmativo, ¿qué tipo de tratamiento se realiza? &                                        &                                          \\
% 			\addlinespace
% 			\hline
% 		\end{tabular}
% 	\end{center}
% \end{table}


\appendixs{Guía de entrevista aplicada a los trabajadores del Vivero Michita}\label{app:guia-entrevista-vivero-michita}

\begin{longtable}{|l|p{6.5cm}|p{3.5cm}|p{3.5cm}|}
	\caption{Guía de entrevista aplicada a los trabajadores del Vivero Michita}\label{tab:entrevista-vivero-michita}                                                                                                                                                \\

	\hline \multicolumn{1}{|c|}{\textbf{N.}} & \multicolumn{1}{c|}{\textbf{Pregunta}}                                                                                         & \multicolumn{1}{c|}{\textbf{Respuesta}} & \multicolumn{1}{c|}{\textbf{Observación}} \\ \hline 
	\endfirsthead

	\multicolumn{4}{c}%
	{{\bfseries \tablename\ \thetable{} -- continuación de la página anterior}}                                                                                                                                                                                     \\
	\hline \multicolumn{1}{|c|}{\textbf{N.}} & \multicolumn{1}{c|}{\textbf{Pregunta}}                                                                                         & \multicolumn{1}{c|}{\textbf{Respuesta}} & \multicolumn{1}{c|}{\textbf{Observación}} \\ \hline
	\endhead

	\hline \multicolumn{4}{|r|}{{Continua en la siguiente página}}                                                                                                                                                                                                  \\ \hline
	\endfoot
	
	\hline \hline
	\endlastfoot

	1                                        & ¿Cuáles son los principales desafíos que enfrentan actualmente en cuanto al consumo de agua en los cultivos del vivero?        &                                         &                                           \\
	2                                        & ¿Qué métodos o herramientas utilizan actualmente para monitorear y controlar el consumo de agua en los cultivos?               &                                         &                                           \\
	3                                        & ¿Qué métodos de riego utilizan actualmente y con qué frecuencia?                                                               &                                         &                                           \\
	4                                        & ¿Cuáles son los principales criterios que tienen en cuenta al decidir cuándo y cuánto regar los cultivos?                      &                                         &                                           \\
	5                                        & ¿Cómo determinan las necesidades de agua de los diferentes tipos de plantas en el vivero?                                      &                                         &                                           \\
	6                                        & ¿Qué problemas o desafíos han experimentado en el pasado con respecto al riego de los cultivos?                                &                                         &                                           \\
	7                                        & ¿Existe algún tratamiento previo al agua antes de su uso en el riego? En caso afirmativo, ¿qué tipo de tratamiento se realiza? &                                         &                                           \\
\end{longtable}



\section{Bibliografía}
\bibliographystyle{biblio/IEEEtran}
\bibliography{biblio/main}
\addcontentsline{toc}{chapter}{\bfseries REFERENCIAS BIBLIOGRÁFICAS}


\end{document}
